\chapter{Architecture globale de la plateforme}

\section{Vue systeme}
La plateforme est organisee autour de quatre blocs techniques :
\begin{enumerate}
\item un frontend de pilotage et de visualisation,
\item une API de service pour les contrats metier,
\item un worker asynchrone pour les traitements longs,
\item une couche de persistance (base relationnelle + stockage d'artefacts).
\end{enumerate}

La Figure~\ref{fig:archi-systeme} illustre cette vue d'ensemble.

\begin{figure}[H]
\centering
\fbox{\parbox{0.92\textwidth}{
\centering
Frontend \(\leftrightarrow\) API \(\leftrightarrow\) Worker\\
API \(\leftrightarrow\) PostgreSQL\\
API et Worker \(\leftrightarrow\) Stockage artefacts (S3/MinIO)
}}
\caption{Architecture technique globale}
\label{fig:archi-systeme}
\end{figure}

La Figure~\ref{fig:archi-systeme} montre que l'API joue le role de facade
contractuelle entre l'interface et les traitements.

\section{Contrats inter-stations}
Le coeur de l'architecture ne se limite pas aux composants techniques ; il
reside dans les contrats de transfert de responsabilite entre stations. Ces
contrats sont resumes dans le Tableau~\ref{tab:contrats-interstations}.

\begin{table}[H]
\centering
\caption{Contrats documentaires inter-stations}
\label{tab:contrats-interstations}
\begin{tabular}{p{3.2cm}p{4cm}p{7.2cm}}
\toprule
\textbf{Flux} & \textbf{Artefact principal} & \textbf{Invariant} \\
\midrule
Data $\rightarrow$ Signal & OHLCV valide & Index temporel ordonne, valeurs exploitables, historique suffisant. \\
Signal $\rightarrow$ Strategie & Findings OOS & Scores interpretables, familles explicites, robustesse tracee. \\
Strategie $\rightarrow$ Backtest & Config WFO-ready & Parametres optimisables identifies, bornes valides, regles coherentes. \\
Backtest $\rightarrow$ Decision & Verdict robuste & WFE, diagnostics de stabilite et tests statistiques exposes. \\
\bottomrule
\end{tabular}
\end{table}

Le Tableau~\ref{tab:contrats-interstations} est reutilise en annexe comme check-list
de verification documentaire.

\section{Choix techniques structurants}
Trois choix ont ete determinants pour la robustesse du systeme :
\begin{itemize}
\item \textbf{API contractuelle explicite} : schemas et endpoints dedies par
station, pour eviter les couplages implicites.
\item \textbf{Worker asynchrone} : separation des traitements longs
(optimisations, batchs, snapshots) et du cycle interactif utilisateur.
\item \textbf{Stockage des artefacts} : conservation des sorties de calcul
pour auditabilite et reproductibilite.
\end{itemize}

\section{Traceabilite et reproductibilite}
La reproductibilite repose sur l'identification du run, la configuration
utilisee, les parametres retenus et les artefacts associes. Cette traque des
preuves est essentielle pour passer d'un outil de demonstration a un outil de
decision.
